\documentclass{article}
\usepackage{../mathnotezh}

\title{Roots of Unity Filter}
\author{tbj}
\date{}

\begin{document}

\maketitle

\section{单位根}

\begin{deff}
    方程$z^n - 1 = 0$的根称为$n$次单位根(Roots of unity).
\end{deff}

易知$n$次单位根有互不相同的$n$个,分别为$\cos \dfrac{2k\pi}{n} + i \sin \dfrac{2k\pi}{n}, k = 0, 1, \cdots, n-1$.

\begin{prop}
    设$\xi = \cos \dfrac{2\pi}{n} + i \sin \dfrac{2\pi}{n}$为$n$次单位根,则

    \[1 + \xi^m + \xi^{2m} + \cdots + \xi^{(n-1)m} =
    \begin{cases}
        n, & n \mid m; \\
        0, & n \nmid m.
    \end{cases}\]
\end{prop}

\begin{proof}
    当$n \mid m$时,有$\xi^m = 1$,故原式$= n$.

    当$n \nmid m$时,有$\xi^m \neq 1$,故原式为等比数列和:
    \[ \text{原式} = \dfrac{1 - (\xi^m)^n}{1 - \xi^m} = 0 \qedhere \]
\end{proof}

\section{Filter}

\begin{thm}
    设$\xi = \cos \dfrac{2\pi}{n} + i \sin \dfrac{2\pi}{n}$为$n$次单位根,多项式$F(x) = a_0 + a_1x + a_2 x^2 + \cdots$(有限多项),
    则\[ a_0 + a_n + a_{2n} + \cdots = \frac{1}{n} \sum_{k=0}^{n-1} F(\xi^k) \]
\end{thm}

\begin{proof}
    易知右端每个$a_m$的系数为

    \[ \frac{1}{n} \sum_{k=0}^{n-1} (\xi^k)^m = 
    \begin{cases}
        1, & n \mid m; \\
        0, & n \nmid m.
    \end{cases} \]

    故右端即为左端.
\end{proof}

这个定理可以帮我们``filter''出那些下标被$n$整除的项的系数和,运用单位根的运算性质.
这个定理就被我们称为``Roots of Unity Filter''.
它也可以用于求解数列中相应项之和,只要运用生成函数(Generating Function)的技术.
有关生成函数的理解可以看这个\href{https://www.bilibili.com/video/BV1QM411A73c/}{视频(BV1QM411A73c)}

笔者没有找到``Roots of Unity Filter''的中文名称,直接翻译可能类似``单位根过滤器'',但中文互联网上没有查找到相应结果.
中文OI(信息学竞赛)社区中有``单位根反演''的说法,指代的算法类似于我们上面所说的定理.
\href{https://notes.sshwy.name/Math/Unit-Root-Inversion/}{一篇文章}给出其英文名为``unit root inversion'',
但英文互联网中也没有相应结果.那我们还是称其为``Roots of Unity Filter''吧.

\begin{cor}
    设$r \in \llbracket 1, n - 1 \rrbracket, G(x) = x^rF(x)$,则

    \[ a_{n - r} + a_{2n - r} + \cdots = \frac{1}{n} \sum_{k=0}^{n-1} G(\xi^k) \]
\end{cor}

这个推论的证明十分简单,我们把它留给读者思考.

\begin{rqq}
    \textbf{用哪个单位根?} \quad 上面我们只使用了$n$个单位根中的一个$\xi = \cos \dfrac{2\pi}{n} + i \sin \dfrac{2\pi}{n}$.
    事实上,对于任何一个与$n$互素的正整数$k$,使用$\xi^k$也可行.但我们仍推荐使用$\xi$.
    我从未见过任何一个问题,其中$\xi^k$可行而$\xi$不可行.

    更深入地讲,当$(n, k) = 1$时,$\{ \xi^k, (\xi^k)^2, \cdots, (\xi^k)^{n - 1} \}$
    为$\{ \xi, \xi^2, \cdots, \xi^{n - 1} \}$的一个排列,因此我们可以说二者没有区别.
\end{rqq}

\section{一些题目}

\begin{exer}
    设$\binom{n}{k}$为二项式系数,求$\binom{2026}{0} + \binom{2026}{3} + \binom{2026}{6} + \cdots + \binom{2026}{2025}$.
\end{exer}

\paragraph{提示} 取$F(x) = (1 + x)^{2026}$,用三次单位根filter.

\begin{exer}
    一个班级中有$a$个男生,$b$个女生.现在要在其中选出若干同学组成一个代表团,要求男生数与女生数之差为4的倍数,共有多少种满足条件的选法?
\end{exer}

\paragraph{解答} 设$F(x) = (1 + x)^a (1 + y)^b$,易知其展开式中$x^my^n$项的系数(记为$c_{mn}$)即为选取$m$个男生与$n$个女生的方法数.
我们要对满足$4 \mid m - n$的$c_{mn}$求和,但$c_{00} = 1$除外.
令$G(x) = F(x, \frac{1}{x})$,则$F(x, y)$中的$x^my^n$项变为$G(x)$中的$x^{m - n}$项.于是对$G(x)$用四次单位根filter即可.

接下来这道题为刚结束的2026年中科大少年班入围考试数学测试倒数第2题.

\begin{exer}
    用3种颜色为$n$边形染色,要求相邻顶点颜色不同,有多少种不同的染色方法?
\end{exer}

\paragraph{解答} 将四种颜色编号为0, 1, 2,则相邻顶点颜色编号之差(模3取余)只能为1, 2.

设n个顶点颜色分别为$s_1, s_2, \cdots, s_n$,相邻顶点颜色之差分别为$d_1, d_2, \cdots, d_n$,
其中$s_{i + 1} \equiv s_i + d_i \Mod{3}, i = 1, 2, \cdots, n - 1, \quad s_1 \equiv s_n + d_n \Mod{4}$.

令$F(x) = (x + x^2)^n$,则其展开式中$x^n$项的系数即为满足$d_1 + d_2 + \cdots + d_n = n$的选法数.
由题意知$3 \mid n$(因为$n$个顶点转一圈回到原处).于是对$F(x)$用3次单位根filter,
这样就得到了满足$3 \mid d_1 + d_2 + \cdots + d_n$的方法数.
再给所得的系数和乘3,因为第一段小绳有3种颜色可选.

\end{document}